% =====================================================================
% Basic MATLAB — Operator
% Penulis : Andreas Syaloom Kurniawan
% Compile : xelatex operator_matlab.tex
% =====================================================================
\documentclass[t,aspectratio=169,11pt]{beamer}
\usetheme[progressbar=frametitle,numbering=fraction]{metropolis}

\usepackage{fontspec}
\setsansfont{DejaVu Sans}
\setmonofont[Scale=0.88]{DejaVu Sans Mono}

\usepackage{listings}
\usepackage{booktabs}
\usepackage{amsmath}
\usepackage{xcolor}

\definecolor{mlorange}{HTML}{D95319}
\definecolor{mlblue}{HTML}{0072BD}
\definecolor{mlgreen}{HTML}{2E7D32}
\definecolor{mlgray}{HTML}{6B6B6B}
\definecolor{codebg}{HTML}{F4F4F1}

\setbeamercolor{alerted text}{fg=mlorange}
\setbeamercolor{progress bar}{fg=mlorange}

\lstset{
  language=Matlab,
  basicstyle=\ttfamily\small,
  keywordstyle=\color{mlblue}\bfseries,
  commentstyle=\color{mlgreen},
  stringstyle=\color{mlorange},
  backgroundcolor=\color{codebg},
  frame=single, rulecolor=\color{codebg},
  framesep=4pt, xleftmargin=4pt, xrightmargin=4pt,
  columns=fullflexible, keepspaces=true,
  showstringspaces=false, upquote=true,
  aboveskip=4pt, belowskip=4pt
}
\newcommand{\cc}[1]{\texttt{\detokenize{#1}}}
\newcommand{\bsl}{\texttt{\textbackslash}}

\title{Basic MATLAB: Operator}
\subtitle{Aritmetika, Relasional, Logika, dan Operator Lain}
\author{Andreas Syaloom Kurniawan}
\institute{Fakultas Ekonomi dan Bisnis, Universitas Gadjah Mada}
\date{Oktober 2026}

\begin{document}

\maketitle

% =====================================================================
\section{Operator Aritmetika}
% =====================================================================

\newcommand{\sub}[1]{{\large\textbf{\textcolor{mlorange}{#1}}}\par\vspace{6pt}}

\begin{frame}{Operator Aritmetika}
\sub{Apa itu operator aritmetika?}
Simbol untuk operasi hitung. Kita mulai dari \alert{angka tunggal} (skalar), misalnya \cc{a = 6} dan \cc{b = 2}.

\vspace{8pt}
\centering
\begin{tabular}{llll}
\toprule
Operator & Arti & Contoh & Hasil \\
\midrule
\cc{+} & tambah  & \cc{a + b} & 8 \\
\cc{-} & kurang  & \cc{a - b} & 4 \\
\cc{*} & kali    & \cc{a * b} & 12 \\
\cc{/} & bagi    & \cc{a / b} & 3 \\
\cc{^} & pangkat & \cc{a ^ b} & 36 \\
\bottomrule
\end{tabular}
\end{frame}

% ---------------------------------------------------------------------
\begin{frame}[fragile]{Operator Aritmetika}
\sub{Mencoba di Command Window}
\begin{columns}[T]
\column{0.50\textwidth}
\begin{lstlisting}
>> a = 6;
>> b = 2;
>> a + b
ans =
     8
>> c = a * b
c =
    12
\end{lstlisting}
\column{0.44\textwidth}
\small
\begin{itemize}
\item \cc{=} menyimpan nilai ke variabel (\cc{a} berisi 6).
\item \cc{;} di akhir baris: hasil \alert{tidak ditampilkan}.
\item Tanpa \cc{;}: hasil langsung tampil.
\item Hasil tanpa nama variabel otomatis disimpan di \cc{ans}.
\end{itemize}
\end{columns}
\end{frame}

% ---------------------------------------------------------------------
\begin{frame}[fragile]{Operator Aritmetika}
\sub{Pembagian: bagi kanan \cc{/} dan bagi kiri \bsl}
MATLAB punya \alert{dua} operator bagi. Bedanya: angka mana yang menjadi pembagi.

\vspace{6pt}
\begin{columns}[T]
\column{0.46\textwidth}
\textbf{\textcolor{mlblue}{Bagi kanan} \cc{/}} (yang biasa)
\begin{lstlisting}
>> 6 / 2
ans =
     3
\end{lstlisting}
\small dibaca: 6 dibagi 2 \quad ($6 \div 2$)
\column{0.46\textwidth}
\textbf{\textcolor{mlorange}{Bagi kiri} \bsl}
\begin{lstlisting}
>> 6 \ 2
ans =
    0.3333
\end{lstlisting}
\small dibaca: 2 dibagi 6 \quad ($2 \div 6$)
\end{columns}

\vspace{8pt}
\small Cara ingat: ujung atas garis miring selalu condong ke \alert{pembagi}.
\cc{6/2} condong ke 2; \cc{6}\bsl\cc{2} condong ke 6.
Untuk angka tunggal pakai \cc{/} saja; kegunaan \bsl{} baru terasa di matriks.
\end{frame}

% ---------------------------------------------------------------------
\begin{frame}[fragile]{Operator Aritmetika}
\sub{Dari angka tunggal ke matriks}
\begin{columns}[T]
\column{0.50\textwidth}
\begin{lstlisting}
>> A = [1 2; 3 4]
A =
     1     2
     3     4
\end{lstlisting}
\small
\begin{itemize}
\item \cc{[ ]} membuat matriks.
\item Spasi (atau koma) memisahkan \alert{kolom}.
\item Titik koma \cc{;} memisahkan \alert{baris}.
\end{itemize}
\column{0.44\textwidth}
Dalam notasi matematika:
\[
A = \begin{bmatrix} 1 & 2 \\ 3 & 4 \end{bmatrix}
\qquad
B = \begin{bmatrix} 5 & 6 \\ 7 & 8 \end{bmatrix}
\]
\small \cc{A} berukuran $2 \times 2$: 2 baris, 2 kolom.\\
\cc{B = [5 6; 7 8]} dibuat dengan cara yang sama.
\end{columns}
\end{frame}

% ---------------------------------------------------------------------
\begin{frame}[fragile]{Operator Aritmetika}
\sub{Tambah dan kurang pada matriks}
Tambah dan kurang selalu dihitung \alert{per posisi} (elemen dengan elemen).

\begin{columns}[T]
\column{0.46\textwidth}
\begin{lstlisting}
>> A + B
ans =
     6     8
    10    12

>> A + 10   % skalar ke semua elemen
ans =
    11    12
    13    14
\end{lstlisting}
\column{0.48\textwidth}
\[
A + B =
\begin{bmatrix} 1+5 & 2+6 \\ 3+7 & 4+8 \end{bmatrix}
=
\begin{bmatrix} 6 & 8 \\ 10 & 12 \end{bmatrix}
\]
\small Syarat: ukuran \cc{A} dan \cc{B} sama, atau salah satunya angka tunggal.
\end{columns}
\end{frame}

% ---------------------------------------------------------------------
\begin{frame}[fragile]{Operator Aritmetika}
\sub{Perkalian: \cc{*} (matriks) vs \cc{.*} (per elemen)}
\begin{columns}[T]
\column{0.48\textwidth}
\textbf{\cc{A * B}: perkalian matriks}\\[2pt]
\small baris \cc{A} $\times$ kolom \cc{B}, lalu dijumlah
\[
\begin{bmatrix} 1{\cdot}5+2{\cdot}7 & 1{\cdot}6+2{\cdot}8 \\ 3{\cdot}5+4{\cdot}7 & 3{\cdot}6+4{\cdot}8 \end{bmatrix}
= \begin{bmatrix} 19 & 22 \\ 43 & 50 \end{bmatrix}
\]
Syarat: jumlah kolom \cc{A} = jumlah baris \cc{B}.
\column{0.48\textwidth}
\textbf{\cc{A .* B}: per elemen}\\[2pt]
\small posisi yang sama dikalikan
\[
\begin{bmatrix} 1{\cdot}5 & 2{\cdot}6 \\ 3{\cdot}7 & 4{\cdot}8 \end{bmatrix}
= \begin{bmatrix} 5 & 12 \\ 21 & 32 \end{bmatrix}
\]
Syarat: ukuran sama (atau salah satunya skalar).
\end{columns}

\vspace{8pt}
\small Titik (\cc{.}) di depan operator artinya \alert{hitung per elemen}. Untuk angka tunggal, \cc{*} dan \cc{.*} hasilnya sama.
\end{frame}

% ---------------------------------------------------------------------
\begin{frame}[fragile]{Operator Aritmetika}
\sub{Pangkat: \cc{^} (matriks) vs \cc{.^} (per elemen)}
\begin{columns}[T]
\column{0.48\textwidth}
\textbf{\cc{A ^ 2}} sama dengan \cc{A * A}
\[
\begin{bmatrix} 1 & 2 \\ 3 & 4 \end{bmatrix}
\begin{bmatrix} 1 & 2 \\ 3 & 4 \end{bmatrix}
= \begin{bmatrix} 7 & 10 \\ 15 & 22 \end{bmatrix}
\]
\small Hanya untuk matriks persegi.
\column{0.48\textwidth}
\textbf{\cc{A .^ 2}} tiap elemen dikuadratkan
\[
\begin{bmatrix} 1^2 & 2^2 \\ 3^2 & 4^2 \end{bmatrix}
= \begin{bmatrix} 1 & 4 \\ 9 & 16 \end{bmatrix}
\]
\small Ini yang biasanya dimaksud saat mengolah data.
\end{columns}
\end{frame}

% ---------------------------------------------------------------------
\begin{frame}[fragile]{Operator Aritmetika}
\sub{Bagi kanan dan bagi kiri pada matriks}
\small Sistem persamaan $2x + y = 3$ dan $x + 3y = 5$ dapat ditulis $A\mathbf{x} = \mathbf{b}$.

\begin{columns}[T]
\column{0.48\textwidth}
\textbf{\textcolor{mlorange}{Bagi kiri}}\\[2pt]\small \cc{x = A}\bsl\cc{b} menyelesaikan \cc{A*x = b}
\begin{lstlisting}
>> A = [2 1; 1 3];
>> b = [3; 5];
>> x = A \ b
x =
    0.8000
    1.4000
\end{lstlisting}
\column{0.48\textwidth}
\textbf{\textcolor{mlblue}{Bagi kanan}}\\[2pt]\small \cc{x = b/A} menyelesaikan \cc{x*A = b}\\
\small (versi baris dari kasus di kiri)
\begin{lstlisting}
>> [3 5] / A
ans =
    0.8000    1.4000
\end{lstlisting}
\small Versi per elemen:\\
\cc{[2 4 6] ./ 2} $\to$ \cc{[1 2 3]}\\
\cc{[2 4 6] .}\bsl\cc{ 2} $\to$ \cc{[1 0.5 0.3333]}
\end{columns}

\vspace{4pt}
\small Pakai \cc{A}\bsl\cc{b}, bukan \cc{inv(A)*b}: lebih cepat dan lebih akurat.
\end{frame}

% ---------------------------------------------------------------------
\begin{frame}[fragile]{Operator Aritmetika}
\sub{Transpose \cc{'}: baris jadi kolom, kolom jadi baris}
\begin{columns}[T]
\column{0.40\textwidth}
\textbf{Vektor $3 \times 1$}
\begin{lstlisting}
>> v = [1; 2; 3];
>> v'
ans =
     1     2     3
\end{lstlisting}
\small ukuran $3 \times 1 \;\to\; 1 \times 3$
\column{0.54\textwidth}
\textbf{Matriks $2 \times 3$}
\begin{lstlisting}
>> M = [1 2 3; 4 5 6];
>> M'
ans =
     1     4
     2     5
     3     6
\end{lstlisting}
\small ukuran $2 \times 3 \;\to\; 3 \times 2$. Cek dengan \cc{size(M')}.
\end{columns}

\vspace{6pt}
\small Catatan: ada juga \cc{.'}. Untuk angka riil hasilnya sama dengan \cc{'}; bedanya hanya pada bilangan kompleks.
\end{frame}

% ---------------------------------------------------------------------
\begin{frame}[fragile]{Operator Aritmetika}
\sub{Urutan prioritas}
\begin{columns}[T]
\column{0.44\textwidth}
\small Dari yang paling dulu dihitung:
\begin{enumerate}
\item kurung \cc{( )}
\item transpose dan pangkat \cc{'} \cc{^} \cc{.^}
\item tanda minus di depan angka
\item \cc{*} \cc{/} \cc{.*} \cc{./} \bsl
\item \cc{+} \cc{-}
\item titik dua \cc{:}
\end{enumerate}
Setingkat: dihitung dari kiri ke kanan.
\column{0.50\textwidth}
\begin{lstlisting}
2 + 3 * 4   % 14, bukan 20
(2 + 3) * 4 % 20

-2^2        % -4  (pangkat dulu)
(-2)^2      %  4

2^3^2       % 64  (kiri ke kanan)
\end{lstlisting}
\small Kalau ragu, \alert{pakai kurung}.
\end{columns}
\end{frame}

% =====================================================================
\section{Operator Relasional}
% =====================================================================

\begin{frame}{Operator Relasional}
\sub{Apa itu operator relasional?}
Operator untuk \alert{membandingkan} dua nilai. Hasil operasi dengan operator relasional
selalu berupa nilai \alert{biner}: \cc{1} artinya \textbf{benar}, \cc{0} artinya \textbf{salah}.

\vspace{8pt}
\centering\small
\begin{tabular}{llll}
\toprule
Operator & Dibaca & Contoh (\cc{a = 5}, \cc{b = 3}) & Hasil \\
\midrule
\cc{==} & sama dengan                 & \cc{a == b} & 0 \\
\cc{~=} & tidak sama dengan           & \cc{a ~= b} & 1 \\
\cc{<}  & kurang dari                 & \cc{a < b}  & 0 \\
\cc{<=} & kurang dari atau sama dengan & \cc{a <= b} & 0 \\
\cc{>}  & lebih dari                  & \cc{a > b}  & 1 \\
\cc{>=} & lebih dari atau sama dengan & \cc{a >= b} & 1 \\
\bottomrule
\end{tabular}
\end{frame}

% ---------------------------------------------------------------------
\begin{frame}[fragile]{Operator Relasional}
\sub{Mencoba di Command Window}
\begin{columns}[T]
\column{0.46\textwidth}
\begin{lstlisting}
>> a = 5;
>> b = 3;
>> a > b
ans =
  logical
   1
>> t = (a == b)
t =
  logical
   0
\end{lstlisting}
\column{0.48\textwidth}
\small
\begin{itemize}
\item 5 lebih dari 3 itu \textbf{benar}, jadi hasilnya \cc{1}.
\item 5 sama dengan 3 itu \textbf{salah}, jadi hasilnya \cc{0}.
\item Hasilnya bertipe \cc{logical}: tipe data khusus yang isinya hanya 1 atau 0. Cek dengan \cc{class(t)}.
\item Nilai logical tetap bisa dihitung: \cc{(a > b) + 1} hasilnya \cc{2}.
\end{itemize}
\end{columns}
\end{frame}

% ---------------------------------------------------------------------
\begin{frame}[fragile]{Operator Relasional}
\sub{Dua hal yang sering tertukar}
\begin{columns}[T]
\column{0.48\textwidth}
\textbf{1. \cc{=} bukan \cc{==}}
\begin{lstlisting}
a = 3     % MENGISI a dengan 3
a == 3    % BERTANYA: apakah a sama
          % dengan 3?  -> 1
\end{lstlisting}
\small Satu tanda sama dengan untuk menyimpan, dua untuk membandingkan.
Menulis \cc{if a = 3} akan error.
\column{0.48\textwidth}
\textbf{2. Tidak sama dengan: \cc{~=}}
\begin{lstlisting}
5 ~= 3    % 1  (benar)
5 != 3    % error di MATLAB
\end{lstlisting}
\small Di Python atau C ditulis \cc{!=}. Di MATLAB tanda ``tidak'' adalah \cc{~} (tilde).
\end{columns}
\end{frame}

% ---------------------------------------------------------------------
\begin{frame}[fragile]{Operator Relasional}
\sub{Membandingkan vektor: dicek per elemen}
Kalau yang dibandingkan vektor, MATLAB mengecek \alert{setiap elemen satu per satu}.
Hasilnya vektor 1/0 dengan ukuran yang sama.

\begin{columns}[T]
\column{0.46\textwidth}
\begin{lstlisting}
>> x = [4 9 1 7];
>> x > 5
ans =
  1x4 logical array
   0   1   0   1
\end{lstlisting}
\column{0.48\textwidth}
\centering\small
\begin{tabular}{lcccc}
\toprule
Posisi & 1 & 2 & 3 & 4 \\
\midrule
\cc{x}        & 4 & 9 & 1 & 7 \\
lebih dari 5? & tidak & ya & tidak & ya \\
\cc{x > 5}    & 0 & \textbf{1} & 0 & \textbf{1} \\
\bottomrule
\end{tabular}
\end{columns}

\vspace{6pt}
\small Dua vektor juga bisa dibandingkan asal ukurannya sama:
\cc{[4 9 1 7] == [4 0 1 0]} $\to$ \cc{[1 0 1 0]}.
\end{frame}

% ---------------------------------------------------------------------
\begin{frame}[fragile]{Operator Relasional}
\sub{Ingat dulu: mengambil elemen berdasarkan posisi}
Selama ini kita mengambil elemen dengan menyebut \alert{nomor posisinya}.

\begin{columns}[T]
\column{0.46\textwidth}
\begin{lstlisting}
>> x = [4 9 1 7];
>> x(2)
ans =
     9
>> x([2 4])
ans =
     9     7
\end{lstlisting}
\column{0.48\textwidth}
\centering\small
\begin{tabular}{lcccc}
\toprule
Posisi & 1 & 2 & 3 & 4 \\
\midrule
\cc{x} & 4 & \textbf{9} & 1 & \textbf{7} \\
\bottomrule
\end{tabular}

\vspace{8pt}
\raggedright
\cc{x([2 4])} artinya: ambil elemen di posisi 2 dan posisi 4.

\vspace{6pt}
Masalahnya: kalau datanya ribuan, kita tidak tahu posisi mana yang nilainya lebih dari 5.
\end{columns}
\end{frame}

% ---------------------------------------------------------------------
\begin{frame}[fragile]{Operator Relasional}
\sub{Logical indexing: mengambil elemen dengan vektor 1/0}
Selain dengan nomor posisi, elemen bisa diambil dengan \alert{vektor logical} yang panjangnya sama:
elemen di posisi bernilai \cc{1} diambil, posisi bernilai \cc{0} dilewati.

\begin{columns}[T]
\column{0.48\textwidth}
\centering\small
\begin{tabular}{lcccc}
\toprule
Posisi & 1 & 2 & 3 & 4 \\
\midrule
\cc{x}    & 4 & 9 & 1 & 7 \\
\cc{idx}  & 0 & 1 & 0 & 1 \\
diambil?  & -- & \textbf{9} & -- & \textbf{7} \\
\bottomrule
\end{tabular}
\column{0.46\textwidth}
\begin{lstlisting}
>> idx = x > 5;
>> x(idx)
ans =
     9     7
\end{lstlisting}
\end{columns}

\vspace{6pt}
\small Inilah yang disebut \alert{logical indexing}: memilih data berdasarkan syarat, bukan berdasarkan posisi.
\end{frame}

% ---------------------------------------------------------------------
\begin{frame}[fragile]{Operator Relasional}
\sub{Logical indexing dalam satu baris}
Dua langkah tadi biasanya ditulis sekaligus:

\begin{columns}[T]
\column{0.46\textwidth}
\begin{lstlisting}
% dua langkah
idx = x > 5;
x(idx)

% satu langkah (lebih umum)
x(x > 5)
\end{lstlisting}
\column{0.48\textwidth}
\small Cara membaca \cc{x(x > 5)}:
\begin{enumerate}
\item hitung dulu yang di dalam kurung: \cc{x > 5} $\to$ \cc{[0 1 0 1]}
\item lalu ambil elemen \cc{x} di posisi yang bernilai 1 $\to$ \cc{[9 7]}
\end{enumerate}
\vspace{4pt}
Dibaca: ``elemen \cc{x} yang lebih dari 5''.
\end{columns}

\vspace{6pt}
\small Catatan: vektornya harus bertipe logical. \cc{x([0 1 0 1])} dengan angka biasa akan error, karena posisi 0 tidak ada.
\end{frame}

% ---------------------------------------------------------------------
\begin{frame}[fragile]{Operator Relasional}
\sub{Kegunaan sehari-hari}
\begin{columns}[T]
\column{0.48\textwidth}
\begin{lstlisting}
x = [4 9 1 7];

x(x > 5)       % ambil nilainya
% [9 7]

find(x > 5)    % di posisi mana?
% [2 4]

x(x > 5) = 0   % ganti nilainya
% [4 0 1 0]
\end{lstlisting}
\column{0.48\textwidth}
\begin{lstlisting}
x = [4 9 1 7];

sum(x > 5)     % ada berapa?
% 2

mean(x > 5)    % berapa proporsinya?
% 0.5000
\end{lstlisting}
\small \cc{sum} dan \cc{mean} bekerja karena 1/0 bisa dijumlahkan: jumlah angka 1 = banyaknya elemen yang memenuhi syarat.
\end{columns}
\end{frame}

% ---------------------------------------------------------------------
\begin{frame}[fragile]{Operator Relasional}
\sub{Jebakan: angka desimal dan teks}
\begin{columns}[T]
\column{0.48\textwidth}
\textbf{Angka desimal}
\begin{lstlisting}
>> 0.1 + 0.2 == 0.3
ans =
  logical
   0
\end{lstlisting}
\small Komputer menyimpan desimal dengan sedikit pembulatan, jadi \cc{0.1 + 0.2} tidak persis 0.3.
Bandingkan dengan toleransi:\\
\cc{abs((0.1+0.2) - 0.3) < 1e-10} $\to$ 1
\column{0.48\textwidth}
\textbf{Teks}
\begin{lstlisting}
>> 'abc' == 'abd'
ans =
  1x3 logical array
   1   1   0
\end{lstlisting}
\small \cc{==} pada teks membandingkan huruf per huruf, dan error kalau panjangnya beda.
Untuk teks pakai \cc{strcmp('abc','abd')} $\to$ 0.
\end{columns}

\vspace{4pt}
\small Membandingkan dua matriks secara utuh: \cc{isequal(A, B)}.
\end{frame}

% =====================================================================
\section{Operator Logika}
% =====================================================================

\begin{frame}{Operator Logika}
\sub{Apa itu operator logika?}
Operator relasional menghasilkan satu syarat (1 atau 0). Operator logika dipakai untuk
\alert{menggabungkan} beberapa syarat atau \alert{membalik} sebuah syarat.
Masukan dan hasilnya sama-sama 1 (benar) atau 0 (salah).

\vspace{6pt}
\centering\footnotesize
\begin{tabular}{lll}
\toprule
Operator & Dibaca & Benar jika \ldots \\
\midrule
\cc{&}    & dan               & kedua syarat benar \\
\cc{|}    & atau              & minimal satu syarat benar \\
\cc{~}    & tidak             & syaratnya salah (dibalik) \\
\cc{&&}   & dan (short-circuit)  & seperti \cc{&}, khusus satu nilai, untuk \cc{if}/\cc{while} \\
\cc{||}   & atau (short-circuit) & seperti \cc{|}, khusus satu nilai, untuk \cc{if}/\cc{while} \\
\cc{xor}  & salah satu saja   & tepat satu syarat benar, tidak boleh dua-duanya \\
\bottomrule
\end{tabular}
\end{frame}

% ---------------------------------------------------------------------
\begin{frame}[fragile]{Operator Logika}
\sub{Kenapa perlu? Contoh: ``apakah x di antara 2 dan 8?''}
\begin{columns}[T]
\column{0.48\textwidth}
\textbf{Cara matematika: \alert{salah} di MATLAB}
\begin{lstlisting}
>> x = 10;
>> 2 < x < 8
ans =
  logical
   1          % padahal 10 > 8 !
\end{lstlisting}
\small MATLAB menghitung dari kiri: \cc{2 < 10} $\to$ 1, lalu \cc{1 < 8} $\to$ 1.
\column{0.48\textwidth}
\textbf{Cara yang benar: dua syarat digabung}
\begin{lstlisting}
>> (x > 2) & (x < 8)
ans =
  logical
   0
\end{lstlisting}
\small Syarat 1: \cc{x > 2} $\to$ 1\\
Syarat 2: \cc{x < 8} $\to$ 0\\
Digabung dengan ``dan'' $\to$ 0. Benar.
\end{columns}
\end{frame}

% ---------------------------------------------------------------------
\begin{frame}[fragile]{Operator Logika}
\sub{\cc{&} (dan): benar jika \emph{kedua} syarat benar}
\small Contoh: mahasiswa \alert{lulus} jika nilai teori $\geq 60$ \textbf{dan} nilai praktik $\geq 60$.

\begin{columns}[T]
\column{0.24\textwidth}
\centering
\begin{tabular}{ccc}
\toprule
A & B & \cc{A & B} \\
\midrule
0 & 0 & 0 \\
0 & 1 & 0 \\
1 & 0 & 0 \\
1 & 1 & \textbf{1} \\
\bottomrule
\end{tabular}
\column{0.36\textwidth}
\begin{lstlisting}
teori   = [70 50 80];
praktik = [65 90 55];

lulus = (teori >= 60) & ...
        (praktik >= 60)
% lulus = [1 0 0]
\end{lstlisting}
\column{0.34\textwidth}
\centering\footnotesize
\begin{tabular}{lccc}
\toprule
Mhs & 1 & 2 & 3 \\
\midrule
teori $\geq$ 60   & 1 & 0 & 1 \\
praktik $\geq$ 60 & 1 & 1 & 0 \\
\cc{&}            & \textbf{1} & 0 & 0 \\
\bottomrule
\end{tabular}
\end{columns}

\vspace{6pt}
\small Satu nilai: \cc{(5 > 3) & (2 > 4)} $\to$ \cc{1 & 0} $\to$ \cc{0}. Bentuk fungsi: \cc{and(A, B)}.
\end{frame}

% ---------------------------------------------------------------------
\begin{frame}[fragile]{Operator Logika}
\sub{\cc{|} (atau): benar jika \emph{minimal satu} syarat benar}
\small Contoh: mahasiswa \alert{remedial} jika nilai teori $< 60$ \textbf{atau} nilai praktik $< 60$.

\begin{columns}[T]
\column{0.24\textwidth}
\centering
\begin{tabular}{ccc}
\toprule
A & B & \cc{A | B} \\
\midrule
0 & 0 & 0 \\
0 & 1 & \textbf{1} \\
1 & 0 & \textbf{1} \\
1 & 1 & \textbf{1} \\
\bottomrule
\end{tabular}
\column{0.36\textwidth}
\begin{lstlisting}
teori   = [70 50 80];
praktik = [65 90 55];

rem = (teori < 60) | ...
      (praktik < 60)
% rem = [0 1 1]
\end{lstlisting}
\column{0.34\textwidth}
\centering\footnotesize
\begin{tabular}{lccc}
\toprule
Mhs & 1 & 2 & 3 \\
\midrule
teori $<$ 60   & 0 & 1 & 0 \\
praktik $<$ 60 & 0 & 0 & 1 \\
\cc{|}         & 0 & \textbf{1} & \textbf{1} \\
\bottomrule
\end{tabular}
\end{columns}

\vspace{6pt}
\small Satu nilai: \cc{(5 > 3) | (2 > 4)} $\to$ \cc{1 | 0} $\to$ \cc{1}. Bentuk fungsi: \cc{or(A, B)}.
\end{frame}

% ---------------------------------------------------------------------
\begin{frame}[fragile]{Operator Logika}
\sub{\cc{~} (tidak): membalik syarat}
\small Benar jadi salah, salah jadi benar. Contoh: siapa yang \alert{tidak lulus}?

\begin{columns}[T]
\column{0.24\textwidth}
\centering
\begin{tabular}{cc}
\toprule
A & \cc{~A} \\
\midrule
0 & \textbf{1} \\
1 & 0 \\
\bottomrule
\end{tabular}
\column{0.40\textwidth}
\begin{lstlisting}
lulus = [1 0 0];
~lulus
% [0 1 1]

~(5 > 3)    % ~1
% 0
\end{lstlisting}
\column{0.30\textwidth}
\small Hasil \cc{~lulus} sama dengan \cc{rem} di slide sebelumnya: ``tidak lulus'' = ``remedial''.
\end{columns}

\vspace{6pt}
\small Tanda \cc{~} yang sama dipakai di \cc{~=} (tidak sama dengan). Bentuk fungsi: \cc{not(A)}.\\
Angka selain 0 dianggap benar: \cc{~5} $\to$ 0, \cc{~0} $\to$ 1.
\end{frame}

% ---------------------------------------------------------------------
\begin{frame}[fragile]{Operator Logika}
\sub{\cc{&&} (dan, short-circuit): idenya}
\small Analogi: satpam gedung mengecek \alert{dulu} apakah tamu membawa KTP. Kalau tidak bawa,
tamu langsung ditolak; satpam \alert{tidak perlu} mencari nama di daftar tamu.

\vspace{4pt}
\begin{columns}[T]
\column{0.48\textwidth}
\textbf{Aturannya}
\begin{itemize}
\item Syarat kiri dicek lebih dulu.
\item Kiri \textbf{salah} $\to$ hasil pasti 0, syarat kanan \alert{dilewati}.
\item Kiri \textbf{benar} $\to$ syarat kanan baru dicek.
\end{itemize}
\column{0.48\textwidth}
\textbf{Bedanya dengan \cc{&}}
\begin{itemize}
\item \cc{&} selalu mengecek \alert{dua-duanya}.
\item \cc{&&} hanya untuk \alert{satu nilai} di tiap sisi, jadi dipakai di \cc{if} dan \cc{while}.
\item Kegunaan utamanya: syarat kiri \alert{melindungi} syarat kanan dari error.
\end{itemize}
\end{columns}
\end{frame}

% ---------------------------------------------------------------------
\begin{frame}[fragile]{Operator Logika}
\sub{\cc{&&}: contoh mengecek nilai mahasiswa ke-\cc{k}}
\lstset{basicstyle=\ttfamily\footnotesize}
\small Data nilai hanya ada untuk 3 mahasiswa. Kita ingin mengecek apakah mahasiswa ke-\cc{k} lulus.

\begin{columns}[T]
\column{0.56\textwidth}
\begin{lstlisting}
nilai = [70 50 80];   % 3 mahasiswa
k = 5;                % mahasiswa ke-5 tidak ada

if k <= numel(nilai) && nilai(k) >= 60
    disp('Lulus')
else
    disp('Tidak lulus / tidak ada')
end
% Tidak lulus / tidak ada
\end{lstlisting}
\small Kalau \cc{&&} diganti \cc{&}: \textcolor{mlorange}{error}, \emph{Index exceeds the number of array elements}.
\column{0.40\textwidth}
\footnotesize
\begin{tabular}{lcc}
\toprule
 & \cc{k = 2} & \cc{k = 5} \\
\midrule
\cc{k <= 3}         & 1 & 0 \\
\cc{nilai(k) >= 60} & 0 & \textit{dilewati} \\
hasil \cc{&&}       & 0 & 0 \\
\bottomrule
\end{tabular}

\vspace{6pt}
\small \cc{numel(nilai)} = banyaknya elemen (3).\\[4pt]
Saat \cc{k = 5}, syarat kiri sudah salah, jadi \cc{nilai(5)} tidak pernah dijalankan.
\end{columns}
\end{frame}

% ---------------------------------------------------------------------
\begin{frame}[fragile]{Operator Logika}
\sub{\cc{||} (atau, short-circuit): idenya}
\small Analogi: mahasiswa boleh ikut ujian jika kehadirannya $\geq$ 75\% \textbf{atau} punya surat izin.
Kalau kehadirannya sudah cukup, dosen \alert{tidak perlu} memeriksa surat izin.

\vspace{4pt}
\begin{columns}[T]
\column{0.48\textwidth}
\textbf{Aturannya}
\begin{itemize}
\item Syarat kiri dicek lebih dulu.
\item Kiri \textbf{benar} $\to$ hasil pasti 1, syarat kanan \alert{dilewati}.
\item Kiri \textbf{salah} $\to$ syarat kanan baru dicek.
\end{itemize}
\column{0.48\textwidth}
\begin{lstlisting}
hadir = 80;      % persen
izin  = false;

if hadir >= 75 || izin
    disp('Boleh ikut ujian')
end
% hadir >= 75 -> 1,
% izin tidak dicek lagi
\end{lstlisting}
\end{columns}
\end{frame}

% ---------------------------------------------------------------------
\begin{frame}[fragile]{Operator Logika}
\sub{\cc{||}: contoh mencari mahasiswa yang perlu ditindaklanjuti}
\small Mahasiswa ke-\cc{k} perlu ditindaklanjuti jika datanya \alert{tidak ada} \textbf{atau} nilainya di bawah 60.

\begin{columns}[T]
\column{0.50\textwidth}
\begin{lstlisting}
nilai = [70 50 80];
k = 5;

if k > numel(nilai) || nilai(k) < 60
    disp('Perlu ditindaklanjuti')
else
    disp('Aman')
end
% Perlu ditindaklanjuti
\end{lstlisting}
\column{0.44\textwidth}
\footnotesize
\begin{tabular}{lccc}
\toprule
 & \cc{k = 1} & \cc{k = 2} & \cc{k = 5} \\
\midrule
\cc{k > 3}         & 0 & 0 & 1 \\
\cc{nilai(k) < 60} & 0 & 1 & \textit{dilewati} \\
hasil \cc{||}      & 0 & 1 & 1 \\
\bottomrule
\end{tabular}

\vspace{6pt}
\small Saat \cc{k = 5}, syarat kiri sudah benar, jadi \cc{nilai(5)} tidak dijalankan. Dengan \cc{|} biasa, baris ini error.
\end{columns}
\end{frame}

% ---------------------------------------------------------------------
\begin{frame}[fragile]{Operator Logika}
\sub{\cc{xor} (salah satu saja)}
\small Benar jika \alert{tepat satu} syarat benar; kalau dua-duanya benar justru salah.
Berupa fungsi, tidak ada simbolnya. Contoh: siapa yang lulus \alert{hanya satu} ujian?

\begin{columns}[T]
\column{0.24\textwidth}
\centering
\begin{tabular}{ccc}
\toprule
A & B & \cc{xor} \\
\midrule
0 & 0 & 0 \\
0 & 1 & \textbf{1} \\
1 & 0 & \textbf{1} \\
1 & 1 & 0 \\
\bottomrule
\end{tabular}
\column{0.36\textwidth}
\begin{lstlisting}
teori   = [70 50 80];
praktik = [65 90 55];

satu = xor(teori >= 60, ...
           praktik >= 60)
% satu = [0 1 1]
\end{lstlisting}
\column{0.34\textwidth}
\centering\footnotesize
\begin{tabular}{lccc}
\toprule
Mhs & 1 & 2 & 3 \\
\midrule
teori $\geq$ 60   & 1 & 0 & 1 \\
praktik $\geq$ 60 & 1 & 1 & 0 \\
\cc{xor}          & 0 & \textbf{1} & \textbf{1} \\
\bottomrule
\end{tabular}
\end{columns}

\vspace{6pt}
\small Bandingkan dengan \cc{|}: mahasiswa 1 lulus dua-duanya, jadi \cc{|} $\to$ 1 tetapi \cc{xor} $\to$ 0.
\end{frame}

% ---------------------------------------------------------------------
\begin{frame}{Operator Logika}
\sub{Ringkasan: kapan pakai yang mana?}
\centering\small
\begin{tabular}{lll}
\toprule
Situasi & Pakai & Contoh \\
\midrule
Menyaring data (vektor/matriks)   & \cc{&} \cc{|} \cc{~} & \cc{x((x > 2) & (x < 8))} \\
Syarat di \cc{if} / \cc{while}    & \cc{&&} \cc{||} \cc{~} & \cc{if n > 0 && rata > 10} \\
Tepat satu syarat yang benar      & \cc{xor} & \cc{xor(A, B)} \\
\bottomrule
\end{tabular}

\vspace{10pt}
\raggedright
Aturan praktis: kalau datanya \alert{lebih dari satu nilai}, pakai \cc{&} dan \cc{|} (satu garis).
Kalau syaratnya \alert{satu nilai} untuk mengambil keputusan, pakai \cc{&&} dan \cc{||} (dua garis).
\end{frame}

% =====================================================================
\section{Operator Lain-lain}
% =====================================================================

\begin{frame}{Operator Lain-lain}
\sub{Ringkasan}
\footnotesize
\centering
\begin{tabular}{lp{0.8\textwidth}}
\toprule
Operator & Fungsi \\
\midrule
\cc{:}   & range (\cc{1:5}, \cc{0:0.1:1}) atau ``semua'' dalam indeks (\cc{A(:,2)}) \\
\cc{( )} & indeks array, argumen fungsi \\
\cc{{ }} & isi cell array atau isi table (\cc{T{1,1}}) \\
\cc{[ ]} & membentuk matriks, menggabungkan (\cc{[A B]} horizontal, \cc{[A; B]} vertikal) \\
\cc{.}   & akses field struct atau kolom table (\cc{S.x}, \cc{T.NamaKolom}) \\
\cc{;}   & pisah baris dalam \cc{[ ]}, atau menekan output di akhir baris \\
\cc{,}   & pisah kolom atau argumen \\
\texttt{\%}   & komentar (\texttt{\%\%} untuk section) \\
\cc{...} & lanjut ke baris berikutnya \\
\bottomrule
\end{tabular}
\end{frame}

% ---------------------------------------------------------------------
\begin{frame}[fragile]{Operator Lain-lain}
\sub{Titik dua \cc{:}}
\begin{columns}[T]
\column{0.48\textwidth}
\textbf{Membuat deret} \cc{awal:langkah:akhir}
\begin{lstlisting}
1:5          % langkah 1 (default)
% [1 2 3 4 5]

0:0.25:1     % langkah 0.25
% [0 0.25 0.5 0.75 1]

10:-2:2      % menurun
% [10 8 6 4 2]
\end{lstlisting}
\column{0.48\textwidth}
\textbf{Artinya ``semua'' dalam indeks}
\begin{lstlisting}
A = [1 2 3; 4 5 6];

A(:, 2)      % semua baris, kolom 2
% [2; 5]

A(1, :)      % baris 1, semua kolom
% [1 2 3]

A(:)         % semua elemen, 1 kolom
% [1; 4; 2; 5; 3; 6]
\end{lstlisting}
\end{columns}
\end{frame}

% ---------------------------------------------------------------------
\begin{frame}[fragile]{Operator Lain-lain}
\sub{Kurung biasa \cc{( )}}
\begin{columns}[T]
\column{0.31\textwidth}
\textbf{Untuk mengindeks}
\begin{lstlisting}
x = [4 9 1 7];
x(3)
% 1

A = [1 2 3; 4 5 6];
A(2, 1)    % baris 2,
% 4        % kolom 1
\end{lstlisting}
\column{0.33\textwidth}
\textbf{Untuk argumen fungsi}
\begin{lstlisting}
max(x)
% 9

round(3.14159, 2)
% 3.14

zeros(2, 3)
% matriks 0 ukuran 2x3
\end{lstlisting}
\column{0.30\textwidth}
\textbf{Untuk urutan hitung}
\begin{lstlisting}
2 + 3 * 4
% 14

(2 + 3) * 4
% 20
\end{lstlisting}
\end{columns}

\vspace{6pt}
\small Indeks di MATLAB mulai dari \alert{1}, bukan 0. Urutannya selalu (baris, kolom).
\end{frame}

% ---------------------------------------------------------------------
\begin{frame}[fragile]{Operator Lain-lain}
\sub{Kurung kurawal \cc{{ }}}
\small Cell array: wadah yang tiap selnya boleh berisi \alert{tipe data berbeda}.

\begin{columns}[T]
\column{0.48\textwidth}
\textbf{Cell array}
\begin{lstlisting}
C = {'Yogyakarta', 2023, [1 2 3]};

C{1}      % ISI sel ke-1
% 'Yogyakarta'

C(1)      % SEL ke-1 (masih cell)
% {'Yogyakarta'}
\end{lstlisting}
\column{0.48\textwidth}
\textbf{Table}
\begin{lstlisting}
T = table([1;2], [3;4], ...
    'VariableNames', {'x','y'});

T{1, 2}   % isi baris 1, kolom 2
% 3

T{:, :}   % seluruh isi jadi matriks
% [1 3; 2 4]
\end{lstlisting}
\end{columns}

\vspace{4pt}
\small Cara ingat: \cc{{ }} mengambil \alert{isi}, \cc{( )} mengambil \alert{wadahnya}.
\end{frame}

% ---------------------------------------------------------------------
\begin{frame}[fragile]{Operator Lain-lain}
\sub{Kurung siku \cc{[ ]}}
\begin{columns}[T]
\column{0.31\textwidth}
\textbf{Membuat vektor/matriks}
\begin{lstlisting}
[1 2 3]      % baris
[1; 2; 3]    % kolom
[1 2; 3 4]   % 2x2
\end{lstlisting}
\column{0.33\textwidth}
\textbf{Menggabungkan matriks}
\begin{lstlisting}
a = [1 2];
b = [3 4];

[a b]     % ke samping
% [1 2 3 4]

[a; b]    % ke bawah
% [1 2; 3 4]
\end{lstlisting}
\column{0.30\textwidth}
\textbf{Menghapus elemen}
\begin{lstlisting}
x = [4 9 1 7];
x(2) = []
% [4 1 7]

y = [];   % kosong
\end{lstlisting}
\end{columns}

\vspace{6pt}
\small Syarat: ke samping jumlah baris sama; ke bawah jumlah kolom sama.
\end{frame}

% ---------------------------------------------------------------------
\begin{frame}[fragile]{Operator Lain-lain}
\sub{Titik \cc{.}}
\small Mengakses bagian bernama di dalam struct atau table.

\begin{columns}[T]
\column{0.48\textwidth}
\textbf{Struct}
\begin{lstlisting}
S.nama = 'Budi';
S.umur = 25;

S.nama
% 'Budi'

S.umur + 1
% 26
\end{lstlisting}
\column{0.48\textwidth}
\textbf{Kolom table}
\begin{lstlisting}
T = table([1;2], [3;4], ...
    'VariableNames', {'x','y'});

T.y          % kolom bernama y
% [3; 4]
\end{lstlisting}
\end{columns}

\vspace{4pt}
\small Jangan tertukar dengan \cc{.*} \cc{./} \cc{.^}: di sana titik artinya ``per elemen''.
\end{frame}

% ---------------------------------------------------------------------
\begin{frame}[fragile]{Operator Lain-lain}
\sub{Titik koma \cc{;}}
\begin{columns}[T]
\column{0.31\textwidth}
\textbf{Memisahkan baris matriks}
\begin{lstlisting}
[1 2; 3 4]
% ans =
%    1   2
%    3   4
\end{lstlisting}
\column{0.33\textwidth}
\textbf{Menyembunyikan hasil}
\begin{lstlisting}
>> a = 5
a =
     5
>> b = 7;
>>            % tidak tampil
\end{lstlisting}
\column{0.30\textwidth}
\textbf{Beberapa perintah satu baris}
\begin{lstlisting}
a = 1; b = 2;
c = a + b;
\end{lstlisting}
\end{columns}

\vspace{6pt}
\small \textbf{Wajib di akhir baris? Tidak.} Pindah baris sudah mengakhiri perintah. Fungsinya hanya menyembunyikan hasil.
\end{frame}

% ---------------------------------------------------------------------
\begin{frame}[fragile]{Operator Lain-lain}
\sub{Titik koma \cc{;} di dalam function}
\small Secara aturan tidak wajib, tetapi di function hampir \alert{selalu} dipakai. Tanpa \cc{;}, hasil antara ikut tercetak setiap kali function dipanggil.

\begin{columns}[T]
\column{0.48\textwidth}
\textbf{File \cc{kuadrat_plus.m}}
\begin{lstlisting}
function y = kuadrat_plus(x)
z = x^2        % lupa titik koma
y = z + 1;
end
\end{lstlisting}
\column{0.48\textwidth}
\textbf{Saat dipanggil}
\begin{lstlisting}
>> a = kuadrat_plus(3);
z =
     9
\end{lstlisting}
\small \cc{z} tetap muncul walau pemanggilnya sudah pakai \cc{;}, karena \cc{;} di luar tidak berlaku untuk baris di dalam function.
\end{columns}

\vspace{4pt}
\small Kebiasaan baik: akhiri setiap baris perhitungan di function dan skrip dengan \cc{;}. Hapus \cc{;} sementara hanya saat ingin melihat hasil antara (debugging).
\end{frame}

% ---------------------------------------------------------------------
\begin{frame}[fragile]{Operator Lain-lain}
\sub{Koma \cc{,}}
\begin{columns}[T]
\column{0.31\textwidth}
\textbf{Memisahkan kolom matriks}
\begin{lstlisting}
[1, 2, 3]
% sama dengan [1 2 3]
\end{lstlisting}
\column{0.33\textwidth}
\textbf{Memisahkan argumen}
\begin{lstlisting}
zeros(2, 3)
round(3.14159, 2)
A(2, 3)
% baris 2, kolom 3
\end{lstlisting}
\column{0.30\textwidth}
\textbf{Memisahkan perintah}
\begin{lstlisting}
>> a = 1, b = 2
a =
     1
b =
     2
\end{lstlisting}
\end{columns}

\vspace{6pt}
\small Bedanya dengan \cc{;}: koma memisahkan tanpa menyembunyikan output.
\end{frame}

% ---------------------------------------------------------------------
\begin{frame}[fragile]{Operator Lain-lain}
\sub{Komentar \texttt{\%} dan section \texttt{\%\%}}
\begin{columns}[T]
\column{0.48\textwidth}
\begin{lstlisting}
%% Bagian 1: Data
x = [4 9 1 7];   % data contoh

%% Bagian 2: Analisis
rata = mean(x);  % rata-rata

%{
Komentar beberapa baris
sekaligus
%}
\end{lstlisting}
\column{0.46\textwidth}
\small
\begin{itemize}
\item Teks setelah \texttt{\%} diabaikan MATLAB: untuk catatan.
\item \texttt{\%\%} di awal baris membagi skrip menjadi \alert{section}. Satu section bisa dijalankan sendiri dengan \textbf{Run Section} (Ctrl+Enter, di Mac Cmd+Enter).
\item \texttt{\%\{} \ldots{} \texttt{\%\}} untuk komentar beberapa baris.
\end{itemize}
\end{columns}
\end{frame}

% ---------------------------------------------------------------------
\begin{frame}[fragile]{Operator Lain-lain}
\sub{Tiga titik \cc{...}}
Menyambung perintah yang terlalu panjang ke baris berikutnya.

\begin{columns}[T]
\column{0.48\textwidth}
\textbf{Rumus panjang}
\begin{lstlisting}
total = harga_1 + harga_2 + ...
        harga_3 + harga_4;
\end{lstlisting}
\textbf{Matriks panjang}
\begin{lstlisting}
x = [1 2 3 4 5 ...
     6 7 8 9 10];
\end{lstlisting}
\column{0.48\textwidth}
\textbf{Argumen fungsi}
\begin{lstlisting}
T = table([1;2], [3;4], ...
    'VariableNames', {'x','y'});
\end{lstlisting}
\small MATLAB membaca semuanya sebagai \alert{satu} perintah. Teks setelah \cc{...} di baris itu dianggap komentar.
\end{columns}
\end{frame}

\end{document}
